% ==============================================================================
% Informe: Agrupació de fenòmens sísmics i detecció d'anomalies (últim mes)
% Compilar amb:  pdflatex informe_sismes.tex  (dues vegades)
% Els valors marcats en vermell [..] són marcadors que s'han d'omplir amb els
% resultats reals obtinguts en executar el codi.
% ==============================================================================
\documentclass[11pt,a4paper]{article}

\usepackage[T1]{fontenc}
\usepackage[utf8]{inputenc}
\usepackage[catalan]{babel}
\usepackage{lmodern}
\usepackage[margin=2.5cm]{geometry}
\usepackage{amsmath,amssymb}
\usepackage{booktabs,array,tabularx}
\usepackage{graphicx}
\usepackage{xcolor}
\usepackage{listings}
\usepackage{algorithm}
\usepackage{algpseudocode}
\usepackage{tikz}
\usetikzlibrary{arrows.meta,positioning,shapes.geometric}
\usepackage{enumitem}
\usepackage{microtype}
\usepackage[colorlinks=true,linkcolor=blue!50!black,citecolor=blue!50!black,urlcolor=blue!50!black]{hyperref}

\floatname{algorithm}{Algorisme}
\renewcommand{\algorithmicrequire}{\textbf{Entrada:}}
\renewcommand{\algorithmicensure}{\textbf{Sortida:}}

% --- Marcadors per als resultats que cal omplir ------------------------------
\newcommand{\pend}[1]{\textcolor{red}{\textbf{[#1]}}}

% --- Figura amb marcador si el fitxer encara no existeix ---------------------
% \figuraopt[amplada]{fitxer}{peu}{etiqueta}
\newcommand{\figuraopt}[4][0.85]{%
  \begin{figure}[htbp]
    \centering
    \IfFileExists{#2}{\includegraphics[width=#1\textwidth]{#2}}{%
      \fbox{\parbox{#1\textwidth}{\centering\vspace{2.2cm}%
      \pend{Inserir figura: \detokenize{#2}}\vspace{2.2cm}}}}
    \caption{#3}\label{#4}
  \end{figure}}

% --- Estil per al codi (només ASCII dins dels llistats) ----------------------
\lstset{
  language=Python,
  basicstyle=\ttfamily\footnotesize,
  keywordstyle=\color{blue!60!black}\bfseries,
  commentstyle=\color{green!40!black}\itshape,
  stringstyle=\color{orange!70!black},
  numbers=left, numberstyle=\tiny\color{gray}, numbersep=6pt,
  breaklines=true, frame=single, rulecolor=\color{gray!50},
  showstringspaces=false, tabsize=4, captionpos=b
}
\renewcommand{\lstlistingname}{Codi}

\newcommand{\R}{\mathbb{R}}
\newcommand{\norm}[1]{\left\lVert #1 \right\rVert}

\title{\textbf{Agrupació de fenòmens sísmics i detecció d'anomalies\\
mitjançant aprenentatge no supervisat}\\[4pt]
\large Anàlisi del catàleg sísmic global de l'últim mes amb \textit{K-means}\\
i puntuacions d'anomalia basades en veïns propers}
\author{\pend{Nom i cognoms de l'autor/a}\\
\pend{Assignatura i curs}\\
\pend{Tòpic de l'exposició: aprenentatge no supervisat / mineria de dades}}
\date{\today}

\begin{document}
\maketitle

\begin{abstract}
\noindent
Aquest informe aborda el problema d'estructurar i filtrar el flux continu de
terratrèmols registrats arreu del món durant l'últim mes. Partint d'un catàleg
públic (magnitud, profunditat, latitud i longitud), s'aplica l'algorisme
\textit{K-means} per classificar els esdeveniments en règims sísmics
diferenciats i s'identifiquen els esdeveniments més excepcionals mitjançant la
distància mitjana als $k$ veïns més propers en l'espai de característiques
normalitzat. Es presenta la fonamentació teòrica de cada tècnica, es compara
amb alternatives (\textit{DBSCAN}, mescles gaussianes, \textit{LOF},
\textit{Isolation Forest}), es descriu la metodologia de manera replicable i es
validen els resultats amb mètriques internes, anàlisi d'estabilitat,
comparació amb una classificació de referència per profunditat i injecció
d'anomalies sintètiques. S'hi inclou també una revisió crítica del codi
Python de partida i la versió millorada (vectoritzada, modular i amb proves
automàtiques). Es discuteixen les limitacions, sobretot la distinció entre
\emph{anòmal estadísticament} i \emph{anòmal físicament}.

\medskip
\noindent\textbf{Paraules clau:} sismologia, \textit{clustering}, K-means,
detecció d'anomalies, normalització \textit{z-score}, aprenentatge no supervisat.
\end{abstract}

\tableofcontents
\newpage

% ==============================================================================
\section{Introducció i plantejament del problema}\label{sec:intro}
% ==============================================================================

\subsection{Context i motivació}
Cada mes, les xarxes sísmiques mundials registren milers de terratrèmols.
La gran majoria són de magnitud petita i passen desapercebuts, però el
conjunt conté informació sobre els règims tectònics actius (zones de subducció,
dorsals oceàniques, falles transformants, sismicitat intraplaca) i, ocasionalment,
esdeveniments excepcionals que requereixen atenció immediata. Revisar manualment
un catàleg d'aquesta mida no és viable, i per això són útils les tècniques
d'aprenentatge no supervisat, que no necessiten etiquetes prèvies:
\emph{(i)} l'agrupació (\textit{clustering}) per descobrir estructura i
\emph{(ii)} la detecció d'anomalies per prioritzar esdeveniments atípics.

Aquest problema mereix ser resolt per tres motius:
\begin{enumerate}[nosep]
  \item \textbf{Interès científic:} resumir un catàleg gran en pocs grups
        interpretables facilita l'exploració i la comparació amb la tectònica
        coneguda.
  \item \textbf{Interès pràctic:} un ordenament d'esdeveniments per grau
        d'excepcionalitat pot servir per prioritzar la revisió humana.
  \item \textbf{Interès docent:} és un cas real, amb dades obertes i
        actualitzades, que il·lustra de manera completa el cicle
        dades $\to$ normalització $\to$ model $\to$ validació, i exposa els
        límits de les hipòtesis geomètriques de \textit{K-means}.
\end{enumerate}

\subsection{Formulació del problema}
Sigui $\mathcal{D}=\{\mathbf{x}_1,\dots,\mathbf{x}_n\}$ el conjunt d'esdeveniments
sísmics registrats en una finestra de 30 dies, on cada esdeveniment
$\mathbf{x}_i\in\R^d$ es descriu per les seves característiques físiques
(magnitud $M$, profunditat $h$ i posició geogràfica). Es volen resoldre dos
subproblemes:
\begin{description}[leftmargin=1.2cm,style=nextline]
  \item[P1 (agrupació)] Trobar una partició $\{C_1,\dots,C_K\}$ de $\mathcal{D}$
        tal que els esdeveniments d'un mateix grup siguin similars entre si i
        diferents dels d'altres grups, i que els grups tinguin una
        interpretació sismològica.
  \item[P2 (anomalies)] Definir una puntuació $s:\mathcal{D}\to\R_{\ge 0}$ tal
        que valors alts identifiquin esdeveniments excepcionals i retornar els
        $m$ amb puntuació més alta (en el codi de partida, $m=5$).
\end{description}

\subsection{Objectius mesurables}\label{sec:objectius}
\begin{table}[htbp]
\centering
\small
\begin{tabularx}{\textwidth}{@{}l>{\raggedright\arraybackslash}X>{\raggedright\arraybackslash}X@{}}
\toprule
\textbf{Id.} & \textbf{Objectiu} & \textbf{Criteri d'acompliment} \\
\midrule
O1 & Carregar i netejar el catàleg dels últims 30 dies. &
     Es documenta el nombre de registres descartats i el motiu; el procés és
     determinista. \\
O2 & Agrupar els esdeveniments amb \textit{K-means} i seleccionar $K$ de
     manera justificada. &
     Silueta mitjana $\bar s\ge 0{,}25$ per al $K$ triat i estabilitat
     $\mathrm{ARI}\ge 0{,}90$ entre 30 llavors aleatòries. \\
O3 & Interpretar físicament els grups. &
     Es reporta l'acord (ARI i NMI) amb la classificació de referència per
     profunditat (superficial/intermèdia/profunda) i es descriu cada grup. \\
O4 & Detectar els $m$ esdeveniments més anòmals. &
     Recuperació $\ge 90\,\%$ d'anomalies sintètiques injectades i
     solapament (Jaccard) amb mètodes alternatius reportat. \\
O5 & Implementació eficient, provada i reproduïble. &
     Totes les proves automàtiques passen; s'informa del temps d'execució
     de la versió \textit{pure Python} i de la vectoritzada per a
     $n=\,$\pend{valor}. \\
\bottomrule
\end{tabularx}
\caption{Objectius i criteris d'acompliment.}\label{tab:objectius}
\end{table}

\subsection{Abast i supòsits}
\textbf{Abast.} Es treballa amb un únic catàleg global (finestra de 30 dies),
amb quatre variables per esdeveniment. No es pretén predir terratrèmols ni
classificar mecanismes focals; només descriure i prioritzar.

\textbf{Supòsits i limitacions assumides:}
\begin{itemize}[nosep]
  \item El catàleg és una mostra representativa però no completa: per sota de
        la \emph{magnitud de completesa} regional els esdeveniments es
        registren de manera desigual.
  \item Els grups esperats són aproximadament compactes en l'espai
        normalitzat (hipòtesi pròpia de \textit{K-means}).
  \item Una finestra d'un mes pot estar dominada per una única seqüència
        sísmica (rèplica d'un gran terratrèmol), cosa que cal tenir present en
        interpretar els resultats.
\end{itemize}

\subsection{Estructura de l'informe}
La secció~\ref{sec:teoria} recull la fonamentació teòrica; la
secció~\ref{sec:metodologia} descriu la metodologia pas a pas;
la secció~\ref{sec:implementacio} detalla la implementació, la revisió del codi
i les proves; la secció~\ref{sec:resultats} presenta els resultats i la
validació; i la secció~\ref{sec:conclusions} recull les conclusions.

% ==============================================================================
\section{Fonamentació teòrica}\label{sec:teoria}
% ==============================================================================

\subsection{Conceptes sismològics bàsics}
\paragraph{Magnitud.} Mesura logarítmica de l'energia alliberada. La magnitud
de moment $M_w$ es defineix a partir del moment sísmic $M_0$ (en N$\cdot$m)
\cite{hanks1979}:
\begin{equation}
  M_w = \tfrac{2}{3}\log_{10} M_0 - 6{,}07 .
\end{equation}
Com que l'escala és logarítmica, un augment d'una unitat de magnitud correspon
a un factor $\approx 32$ en energia, de manera que la distribució de
magnituds és molt asimètrica.

\paragraph{Llei de Gutenberg--Richter.} El nombre d'esdeveniments $N$ amb
magnitud $\ge M$ segueix \cite{gutenberg1944}:
\begin{equation}
  \log_{10} N(\ge M) = a - bM,
\end{equation}
amb $b\approx 1$ a escala global. Això implica que els grans terratrèmols són
rars per naturalesa; en aquest treball són, doncs, candidats naturals a ser
\emph{anomalies estadístiques}.

\paragraph{Profunditat i règims tectònics.} La profunditat focal es clasifica
convencionalment en sismes \emph{superficials} ($h<70$~km), \emph{intermedis}
($70\le h\le 300$~km) i \emph{profunds} ($h>300$~km). Els intermedis i profunds
s'associen gairebé exclusivament a zones de subducció, on la placa que
s'enfonsa genera la zona de Wadati--Benioff \cite{benioff1949}. La major part
de la sismicitat superficial s'agrupa al llarg dels límits de placa
\cite{bird2003}. Aquesta classificació serveix com a \emph{referència externa}
per validar els grups no supervisats (secció~\ref{sec:validacio}).

\subsection{Preparació de les dades: normalització \textit{z-score}}
Les variables tenen unitats i escales molt diferents (magnitud $\sim 1$--$8$,
profunditat $0$--$700$~km, longitud $-180$--$180^\circ$). Com que
\textit{K-means} i les distàncies de veïns són sensibles a l'escala, es
transformen les variables:
\begin{equation}
  z_{ij} = \frac{x_{ij}-\mu_j}{\sigma_j},\qquad
  \mu_j=\frac1n\sum_{i=1}^n x_{ij},\quad
  \sigma_j=\sqrt{\frac1n\sum_{i=1}^n (x_{ij}-\mu_j)^2}.
\end{equation}
Si una columna és constant ($\sigma_j=0$) es fixa $\sigma_j=1$ per evitar la
divisió per zero. Després de la transformació cada variable té mitjana 0 i
desviació típica 1, de manera que cap domina la distància només per les seves
unitats. Alternatives raonables són la normalització \textit{min--max}
(molt sensible a valors extrems, justament els que busquem) i l'escalat robust
amb mediana i IQR; es descarta el primer per aquesta raó i es manté
\textit{z-score} per simplicitat i pel seu ús estàndard amb \textit{K-means}.

\subsection{Representació de la posició geogràfica}\label{sec:geo}
Utilitzar latitud i longitud com a coordenades euclidianes té dos problemes:
(i) la longitud és periòdica (170$^\circ$E i 170$^\circ$W són propers, però
difereixen en 340$^\circ$), i (ii) un grau de longitud correspon a una
distància que depèn de la latitud. Una solució estàndard és projectar cada
punt sobre l'esfera:
\begin{equation}
  \mathbf{p}=R\begin{pmatrix}\cos\varphi\cos\lambda\\ \cos\varphi\sin\lambda\\ \sin\varphi\end{pmatrix},\quad R=6371~\text{km},
\end{equation}
on $\varphi$ és la latitud i $\lambda$ la longitud. La distància euclidiana
(corda) entre dos punts d'aquest tipus aproxima molt bé la distància de
cercle màxim, que té l'expressió \emph{haversine} \cite{sinnott1984}:
\begin{equation}
  d = 2R\arcsin\sqrt{\sin^2\!\tfrac{\Delta\varphi}{2}
      +\cos\varphi_1\cos\varphi_2\sin^2\!\tfrac{\Delta\lambda}{2}}.
\end{equation}
En l'informe es comparen dues representacions de les característiques:
\textbf{(A)} $(M,h,\varphi,\lambda)$, com al codi de partida, i
\textbf{(B)} $(M,h,x,y,z)$ amb coordenades cartesianes sobre l'esfera.

\subsection{Algorisme \textit{K-means}}
\textit{K-means} \cite{macqueen1967,lloyd1982} cerca $K$ centroides
$\boldsymbol\mu_1,\dots,\boldsymbol\mu_K$ que minimitzen la suma de quadrats
intragrup (inèrcia):
\begin{equation}
  J(C,\boldsymbol\mu)=\sum_{k=1}^{K}\sum_{\mathbf{z}_i\in C_k}\norm{\mathbf{z}_i-\boldsymbol\mu_k}^2 .
\end{equation}
L'algorisme de Lloyd alterna dos passos fins a la convergència:
\begin{align}
  \text{Assignació:}\quad & c_i=\arg\min_{k}\norm{\mathbf{z}_i-\boldsymbol\mu_k}^2,\\
  \text{Actualització:}\quad & \boldsymbol\mu_k=\frac{1}{|C_k|}\sum_{\mathbf{z}_i\in C_k}\mathbf{z}_i .
\end{align}
Cada iteració no augmenta $J$, de manera que l'algorisme convergeix a un
mínim \emph{local}. El cost per iteració és $O(nKd)$.

\paragraph{Limitacions rellevants per a aquest problema.}
\begin{itemize}[nosep]
  \item Depèn de la inicialització: el codi de partida tria $K$ punts a
        l'atzar; \textit{k-means++} \cite{arthur2007} redueix aquest risc
        ponderant per distància.
  \item Suposa grups esfèrics i de mida semblant; falla amb formes allargades
        com una fossa de subducció.
  \item Cal fixar $K$ a priori (secció~\ref{sec:triaK}).
  \item És sensible als valors extrems, que arrosseguen els centroides: això
        és rellevant perquè els esdeveniments anòmals formen part de les dades.
\end{itemize}

\subsection{Selecció del nombre de grups $K$}\label{sec:triaK}
\paragraph{Mètode del colze.} Es representa $J(K)$ i es busca el punt on la
millora marginal es redueix de manera notable. És subjectiu, per això es
complementa amb mètriques numèriques.

\paragraph{Coeficient de silueta} \cite{rousseeuw1987}. Per a cada punt $i$, sigui
$a(i)$ la distància mitjana als punts del seu grup i $b(i)$ la menor distància
mitjana a un altre grup:
\begin{equation}
  s(i)=\frac{b(i)-a(i)}{\max\{a(i),b(i)\}}\in[-1,1].
\end{equation}
Valors propers a 1 indiquen una assignació clara. Com a guia, valors per sota
de $0{,}25$ indiquen que no hi ha una estructura substancial
\cite{kaufman1990}.

\paragraph{Índex de Davies--Bouldin} \cite{davies1979}. Mitjana, per a cada
grup, de la pitjor relació entre dispersió intragrup i separació entre
centroides; valors més baixos són millors:
\begin{equation}
  \mathrm{DB}=\frac1K\sum_{k=1}^K\max_{l\neq k}\frac{S_k+S_l}{\norm{\boldsymbol\mu_k-\boldsymbol\mu_l}},
\end{equation}
amb $S_k$ la dispersió mitjana del grup $k$ respecte del seu centroide.

\paragraph{Estabilitat i comparació de particions.} Per avaluar si la partició
depèn de la llavor es calcula l'\emph{Adjusted Rand Index} (ARI)
\cite{hubert1985} entre execucions; val 1 per a particions idèntiques i
$\approx 0$ per a particions aleatòries. La \emph{informació mútua
normalitzada} (NMI) s'usa de manera complementària per comparar amb la
classificació de referència.

\subsection{Detecció d'anomalies}
Una anomalia és una observació que s'allunya de manera notable del
comportament habitual de la majoria \cite{chandola2009}. En aprenentatge no
supervisat no hi ha etiquetes; s'assigna una puntuació i es ressalten les més
altes.

\paragraph{Distància mitjana als $k$ veïns (mètode del codi).} Sigui
$\mathcal{N}_k(\mathbf{z})$ el conjunt dels $k$ veïns més propers de
$\mathbf{z}$ (excloent-lo a ell mateix). La puntuació és
\cite{ramaswamy2000}
\begin{equation}
  s_k(\mathbf{z})=\frac1k\sum_{\mathbf{y}\in\mathcal{N}_k(\mathbf{z})}\norm{\mathbf{z}-\mathbf{y}} .
\end{equation}
És intuïtiva i no paramètrica, però és una mesura \emph{global}: no té en
compte que algunes regions són naturalment més denses que d'altres. El cost
ingenu és $O(n^2d)$.

\paragraph{\textit{Local Outlier Factor} (LOF).} \cite{breunig2000} Compara la
densitat local d'un punt amb la dels seus veïns. Amb la distància d'accessibilitat
$\mathrm{rd}_k(\mathbf{z},\mathbf{y})=\max\{d_k(\mathbf{y}),\,d(\mathbf{z},\mathbf{y})\}$,
on $d_k(\mathbf{y})$ és la distància de $\mathbf{y}$ al seu $k$-èsim veí:
\begin{equation}
  \mathrm{lrd}_k(\mathbf{z})=\Bigg(\frac{1}{k}\sum_{\mathbf{y}\in\mathcal{N}_k(\mathbf{z})}\mathrm{rd}_k(\mathbf{z},\mathbf{y})\Bigg)^{-1},\quad
  \mathrm{LOF}_k(\mathbf{z})=\frac1k\sum_{\mathbf{y}\in\mathcal{N}_k(\mathbf{z})}\frac{\mathrm{lrd}_k(\mathbf{y})}{\mathrm{lrd}_k(\mathbf{z})}.
\end{equation}
Valors $\gg 1$ indiquen que el punt és molt menys dens que els seus veïns.

\paragraph{\textit{Isolation Forest}.} \cite{liu2008} Construeix arbres
aleatoris que divideixen l'espai; els punts anòmals s'aïllen amb poques
particions. La puntuació és
$s(\mathbf{z},n)=2^{-\mathbb{E}[h(\mathbf{z})]/c(n)}$, on $h$ és la longitud del
camí i $c(n)=2H(n-1)-2(n-1)/n$ un factor de normalització ($H$ és el nombre
harmònic).

\subsection{Comparació amb enfocaments existents}
\begin{table}[htbp]
\centering
\small
\begin{tabularx}{\textwidth}{@{}>{\raggedright\arraybackslash}p{2.7cm}>{\raggedright\arraybackslash}X>{\raggedright\arraybackslash}X@{}}
\toprule
\textbf{Tècnica} & \textbf{Avantatges} & \textbf{Inconvenients en aquest problema} \\
\midrule
\textit{K-means} (escollit) & Simple, ràpid, interpretable via centroides;
  base de referència habitual. & Grups esfèrics; $K$ fix; sensible a extrems. \\
\textit{DBSCAN} \cite{ester1996} & Detecta formes arbitràries i marca soroll
  (útil per a falles i fosses). & Sensible a $\varepsilon$ i \textit{minPts};
  densitats molt desiguals entre regions. \\
Mescles gaussianes (GMM) & Grups el·líptics i assignació probabilística. &
  Més paràmetres; pot ser inestable amb pocs punts per grup. \\
Agrupació espai-temps \cite{zaliapin2013} & Separa seqüències de rèpliques de
  la sismicitat de fons amb criteris físics. & Requereix el temps d'origen i
  més modelatge; fora de l'abast d'aquest treball (línia de millora). \\
\midrule
Distància als $k$ veïns (escollit) & Senzilla, sense hipòtesi de distribució. &
  No local; $O(n^2)$ sense estructures espacials. \\
LOF & Sensible a la densitat local. & Més cara; depèn de $k$. \\
\textit{Isolation Forest} & Escalable, $\approx O(n\log n)$. &
  Aleatòria; menys interpretable. \\
\bottomrule
\end{tabularx}
\caption{Comparació de les tècniques escollides amb alternatives.}\label{tab:comparacio}
\end{table}

L'ús de \textit{K-means} en dades sísmiques és habitual com a primera
aproximació exploratòria, però la bibliografia especialitzada recomana
mètodes que incorporin el temps i la física dels processos de ruptura
\cite{zaliapin2013}. Per això, aquest informe tracta \textit{K-means} com una
\emph{línia base} i avalua amb rigor quan els seus resultats són
interpretables i quan no.

% ==============================================================================
\section{Metodologia i solució proposada}\label{sec:metodologia}
% ==============================================================================

\subsection{Visió general}
La figura~\ref{fig:pipeline} resumeix el procés complet. Cada etapa es descriu
a continuació amb prou detall perquè una altra persona pugui replicar-lo.

\begin{figure}[htbp]
\centering
\begin{tikzpicture}[
  caixa/.style={draw, rounded corners, align=center, minimum width=2.7cm,
                minimum height=1.0cm, fill=blue!6, font=\small},
  >=Stealth]
  \node[caixa] (a) at (0,0)    {Catàleg sísmic\\(CSV, 30 dies)};
  \node[caixa] (b) at (3.6,0)  {Neteja i\\validació};
  \node[caixa] (c) at (7.2,0)  {Característiques\\+ \textit{z-score}};
  \node[caixa, fill=green!8] (d) at (3.6,-2.2) {\textit{K-means}\\(P1: agrupació)};
  \node[caixa, fill=orange!10] (e) at (10.0,-2.2) {Puntuació $s_k$\\(P2: anomalies)};
  \node[caixa, fill=gray!12] (f) at (6.8,-4.4) {Validació i\\interpretació};
  \draw[->] (a)--(b); \draw[->] (b)--(c);
  \draw[->] (c)--(d); \draw[->] (c)--(e);
  \draw[->] (d)--(f); \draw[->] (e)--(f);
\end{tikzpicture}
\caption{Diagrama del procés: de les dades en brut a la interpretació.}
\label{fig:pipeline}
\end{figure}

\subsection{Pas 1: Obtenció de les dades}
Les dades provenen del catàleg d'esdeveniments de l'USGS \cite{usgs}, que és
públic i s'actualitza de manera contínua. Es descarrega un CSV amb els
esdeveniments entre \pend{data inicial} i \pend{data final} (30 dies) i es desa
com a \texttt{earthquakes.csv} al directori del projecte, de manera que
l'anàlisi no depèn de la xarxa en executar-se. El fitxer emprat per a
l'informe s'inclou al repositori com a \emph{dades d'exemple}
(secció~\ref{sec:reproducibilitat}), ja que el catàleg canvia amb el temps i,
altrament, els resultats no es podrien reproduir exactament. Camps
utilitzats: \texttt{mag}, \texttt{depth}, \texttt{latitude},
\texttt{longitude} i \texttt{place} (etiqueta de l'esdeveniment).

\subsection{Pas 2: Neteja i validació}
Es descarten les files amb valors no numèrics, buits o no finits (\texttt{NaN},
$\pm\infty$) en alguna de les quatre variables i les latituds o longituds fora
de $[-90,90]$ i $[-180,180]$. Es registra el recompte de files descartades
(objectiu O1), que s'informa a la taula~\ref{tab:dades}. Una diferència
deliberada respecte al codi de partida és que \emph{no} s'imputa una
profunditat de 10~km quan falta, perquè això introduiria un valor inventat que
podria agrupar-se artificialment.

\subsection{Pas 3: Característiques i normalització}
Es construeixen les dues representacions (A) i (B) de la
secció~\ref{sec:geo} i s'apliquen \textit{z-score} columna a columna. Els
paràmetres $\mu_j,\sigma_j$ es conserven per poder interpretar els centroides
en unitats originals.

\subsection{Pas 4: Agrupació amb \textit{K-means}}
\begin{algorithm}[htbp]
\caption{\textit{K-means} amb inicialització aleatòria i criteri de parada}
\label{alg:kmeans}
\begin{algorithmic}[1]
\Require $\mathbf{Z}\in\R^{n\times d}$, $K$, màxim d'iteracions $T$, tolerància $\tau$, llavor $s$
\Ensure etiquetes $c\in\{1..K\}^n$, centroides $\boldsymbol\mu$, inèrcia $J$
\State \textbf{si} $K<1$ \textbf{o} $K>n$ \textbf{llavors} llançar error
\State Inicialitzar el generador aleatori amb la llavor $s$
\State $\boldsymbol\mu\gets$ $K$ files de $\mathbf{Z}$ triades sense reemplaçament
\For{$t=1,\dots,T$}
  \State $c_i\gets\arg\min_k\norm{\mathbf{z}_i-\boldsymbol\mu_k}^2$ per a tot $i$
  \State $\boldsymbol\mu'_k\gets$ mitjana dels punts amb $c_i=k$ \quad (si el grup és buit, conservar $\boldsymbol\mu_k$)
  \If{$\norm{\boldsymbol\mu'-\boldsymbol\mu}<\tau$} \State $\boldsymbol\mu\gets\boldsymbol\mu'$; \textbf{sortir} \EndIf
  \State $\boldsymbol\mu\gets\boldsymbol\mu'$
\EndFor
\State \Return $c,\boldsymbol\mu,J(c,\boldsymbol\mu)$
\end{algorithmic}
\end{algorithm}

\paragraph{Decisions i justificació.} La taula~\ref{tab:decisions} recull les
decisions principals, les alternatives considerades i el motiu de l'elecció.

\begin{table}[htbp]
\centering
\small
\begin{tabularx}{\textwidth}{@{}>{\raggedright\arraybackslash}p{3.0cm}>{\raggedright\arraybackslash}p{3.5cm}>{\raggedright\arraybackslash}X@{}}
\toprule
\textbf{Decisió} & \textbf{Alternatives} & \textbf{Motiu} \\
\midrule
Escalat \textit{z-score} & \textit{min--max}, escalat robust &
  \textit{Min--max} comprimeix tot el rang segons els extrems que volem
  detectar; \textit{z-score} és l'estàndard amb \textit{K-means}. \\
Posició: (A) lat/lon i (B) esfera 3D & Només (A) &
  (B) evita el salt de longitud i les distorsions per latitud; (A) es manté
  per coherència amb el codi de partida i per quantificar la diferència. \\
$K$ per silueta + DB + colze & Fixar $K=3$ a priori &
  $K=3$ es motiva per la classificació de profunditat, però cal comprovar que
  les dades ho suporten; es reporta tant $K=3$ com el $K$ òptim empíric. \\
Llavor fixa i 30 repeticions & Una sola execució &
  La llavor fixa assegura reproducibilitat; les 30 repeticions mesuren
  l'estabilitat (ARI). \\
$k=5$ veïns per a anomalies & $k\in\{3,5,10,20\}$ &
  $k=5$ és el valor del codi de partida; es mostra la sensibilitat a $k$ i
  s'avalua amb \textit{LOF} i \textit{Isolation Forest}. \\
$m=5$ anomalies principals & Llindar per percentil &
  Un recompte fix és fàcil d'interpretar en un informe; es compara també amb
  el percentil 99. \\
\bottomrule
\end{tabularx}
\caption{Decisions de disseny, alternatives i justificació.}\label{tab:decisions}
\end{table}

\subsection{Pas 5: Detecció d'anomalies}
Per a cada punt es calcula $s_k$ amb $k=5$ i s'ordena de manera descendent. Els
$m$ primers es consideren anomalies. Per no confondre magnitud amb
excepcionalitat, s'analitza quin component (magnitud, profunditat o posició)
explica cada puntuació alta, comparant el valor $z$ de cada variable.

\subsection{Pas 6: Validació}\label{sec:metodo-validacio}
\begin{enumerate}[nosep]
  \item \textbf{Mètriques internes:} inèrcia, silueta i Davies--Bouldin per a
        $K=2,\dots,10$.
  \item \textbf{Estabilitat:} ARI entre les particions de 30 llavors diferents.
  \item \textbf{Referència externa:} taula de contingència i ARI/NMI entre els
        grups i la classificació per profunditat (70 i 300~km).
  \item \textbf{Anomalies sintètiques:} s'injecten $m_s=\,$\pend{valor} esdeveniments
        artificials extrems (p.~ex. $M\ge 8{,}5$ a $h>600$~km o
        posicions aïllades) i es mesura la recuperació en els $m_s$ primers
        llocs.
  \item \textbf{Comparació de mètodes:} solapament (Jaccard) dels $m$ primers
        entre $s_k$, \textit{LOF} i \textit{Isolation Forest}.
  \item \textbf{Contrast amb fonts externes:} es comproven els anòmals
        trobats amb el canal d'esdeveniments significatius de l'USGS
        \cite{usgs}.
\end{enumerate}

\subsection{Complexitat}
Per a $n$ esdeveniments, $d$ característiques, $K$ grups i $T$ iteracions:
\textit{K-means} costa $O(TnKd)$, el càlcul ingenu de $s_k$ costa
$O(n^2d)$ en temps i $O(n)$ en memòria si s'ordena fila a fila; amb un arbre
KD el cost mitjà es redueix a $\approx O(n\log n)$ en dimensió baixa.

% ==============================================================================
\section{Implementació, qualitat del codi i reproducibilitat}\label{sec:implementacio}
% ==============================================================================

\subsection{Revisió del codi de partida}
El \textit{script} de partida (\texttt{presentation.py}) implementa tot el
procés amb Python pur i \texttt{matplotlib}. La revisió ha detectat els
punts següents, que la versió millorada resol:

\begin{table}[htbp]
\centering
\small
\begin{tabularx}{\textwidth}{@{}>{\raggedright\arraybackslash}p{4.4cm}>{\raggedright\arraybackslash}X@{}}
\toprule
\textbf{Observació} & \textbf{Correcció aplicada} \\
\midrule
\textbf{Error en el gràfic:} es fa
\texttt{[data\_raw[i][0] for i in distancies\_veins]}, però els elements de
\texttt{distancies\_veins} són parells \texttt{(distància, índex)}; això
provoca un \texttt{TypeError}. A més, no es limita als 5 primers. &
Iterar sobre \texttt{distancies\_veins[:5]} i desempaquetar l'índex
(\texttt{for \_, i in ...}). \\
\texttt{import matplotlib} duplicat dins el bloc de lectura. & Un únic
\textit{import} a l'inici. \\
Bucles en Python pur; la detecció d'anomalies és $O(n^2)$ amb un bucle
intern que crea i ordena una llista per punt. &
Vectorització amb NumPy i \texttt{NearestNeighbors} (arbre KD) de
\textit{scikit-learn} \cite{pedregosa2011}. \\
Tot el codi és un \textit{script} lineal amb valors fixos ($K=3$, 25
iteracions, $k=5$). & Funcions amb responsabilitat única i paràmetres
explícits. \\
Profunditat per defecte de 10~km si falta; entrades errònies s'ignoren en
silenci. & Es descarten i es comptabilitzen les files invàlides. \\
Latitud i longitud tractades com a coordenades euclidianes. &
Variant (B) amb coordenades sobre l'esfera. \\
Sense proves, \textit{README} ni \texttt{requirements.txt}. &
Proves amb \texttt{pytest}, \textit{README} i dependències fixades. \\
Els centroides inicials es trien a l'atzar. &
S'afegeix una llavor i l'opció \textit{k-means++} per comparar. \\
\bottomrule
\end{tabularx}
\caption{Problemes del codi de partida i correccions.}\label{tab:revisio}
\end{table}

\subsection{Arquitectura modular}
S'organitza el projecte en mòduls amb una responsabilitat clara, separant
dades, lògica i presentació:
\begin{lstlisting}[language={},numbers=none,caption={Estructura del projecte.}]
sismes/
  README.md
  requirements.txt
  earthquakes.csv          # dades d'exemple (finestra de 30 dies)
  sismes/
    __init__.py
    dades.py               # lectura i validacio del CSV
    preprocessament.py     # z-score, coordenades esferiques
    clustering.py          # K-means i metriques de validacio
    anomalies.py           # puntuacio kNN (i comparacio amb LOF/IF)
    visualitzacio.py       # grafics (Matplotlib)
  main.py                  # orquestra el proces i imprimeix resultats
  tests/
    test_sismes.py
\end{lstlisting}

\subsection{Funcions principals}
Els codis~\ref{lst:dades}--\ref{lst:anomalies} mostren les funcions
essencials, totes amb \textit{docstring}. Es fan servir noms descriptius i els
comentaris expliquen el \emph{perquè} de cada decisió.

\begin{lstlisting}[caption={Lectura i validacio de dades (\texttt{dades.py}).},label={lst:dades}]
import csv
import math
import numpy as np

# Alias acceptats per a cada columna segons la procedencia del CSV.
ALIAS = {
    "mag": ("magnitude", "mag", "Mag"),
    "profunditat": ("depth", "Depth", "depth_km"),
    "lat": ("latitude", "lat", "Lat"),
    "lon": ("longitude", "lon", "Lon"),
}


def _valor(fila, noms):
    """Retorna el primer camp present en 'fila' convertit a float."""
    for nom in noms:
        if nom in fila:
            return float(fila[nom])
    raise KeyError(noms)


def carregar_csv(ruta):
    """Llegeix el cataleg i retorna (X, identificadors).

    X te forma (n, 4) amb columnes [mag, profunditat, lat, lon].
    Les files incompletes o no finites es descarten (no s'inventen valors).
    Llanca FileNotFoundError si la ruta no existeix i ValueError si no
    queda cap fila valida.
    """
    ids, files = [], []
    with open(ruta, encoding="utf-8", newline="") as f:
        for idx, fila in enumerate(csv.DictReader(f), start=1):
            try:
                valors = [_valor(fila, noms) for noms in ALIAS.values()]
            except (ValueError, KeyError):
                continue
            lat, lon = valors[2], valors[3]
            if not all(math.isfinite(v) for v in valors):
                continue
            if not (-90 <= lat <= 90 and -180 <= lon <= 180):
                continue
            ids.append(fila.get("place") or fila.get("id") or f"Sisme #{idx}")
            files.append(valors)
    if not files:
        raise ValueError("Cap fila valida al fitxer: " + ruta)
    return np.array(files, dtype=float), ids
\end{lstlisting}

\begin{lstlisting}[caption={Preprocessament (\texttt{preprocessament.py}).},label={lst:prep}]
import numpy as np


def estandarditzar(X):
    """Aplica z-score per columnes; retorna (Z, mitjanes, desviacions).

    Si una columna es constant, la desviacio es fixa a 1 per evitar
    dividir per zero.
    """
    mitjanes = X.mean(axis=0)
    desviacions = X.std(axis=0)          # poblacional, com a la memoria
    desviacions[desviacions == 0] = 1.0
    return (X - mitjanes) / desviacions, mitjanes, desviacions


def cartesianes_esfera(lat, lon, radi=6371.0):
    """Passa (lat, lon) en graus a coordenades 3D (km) sobre l'esfera.

    Evita la discontinuitat de la longitud a +-180 graus.
    """
    phi, lam = np.radians(lat), np.radians(lon)
    return radi * np.column_stack([
        np.cos(phi) * np.cos(lam),
        np.cos(phi) * np.sin(lam),
        np.sin(phi),
    ])
\end{lstlisting}

\begin{lstlisting}[caption={\textit{K-means} vectoritzat (\texttt{clustering.py}).},label={lst:kmeans}]
import numpy as np


def kmeans(X, k, n_iter=100, tol=1e-6, llavor=42):
    """K-means (Lloyd). Retorna (etiquetes, centroides, inercia).

    Llanca ValueError si k no esta entre 1 i el nombre de punts.
    """
    if not 1 <= k <= len(X):
        raise ValueError("k ha d'estar entre 1 i n")
    rng = np.random.default_rng(llavor)
    centres = X[rng.choice(len(X), size=k, replace=False)]
    for _ in range(n_iter):
        # distancies al quadrat punt-centroide: forma (n, k)
        d2 = ((X[:, None, :] - centres[None, :, :]) ** 2).sum(axis=2)
        etiquetes = d2.argmin(axis=1)
        nous = np.array([
            X[etiquetes == j].mean(axis=0) if np.any(etiquetes == j)
            else centres[j]              # grup buit: es conserva el centroide
            for j in range(k)
        ])
        convergit = np.linalg.norm(nous - centres) < tol
        centres = nous
        if convergit:
            break
    d2 = ((X[:, None, :] - centres[None, :, :]) ** 2).sum(axis=2)
    etiquetes = d2.argmin(axis=1)
    inercia = d2[np.arange(len(X)), etiquetes].sum()
    return etiquetes, centres, inercia
\end{lstlisting}

\begin{lstlisting}[caption={Puntuació d'anomalia (\texttt{anomalies.py}).},label={lst:anomalies}]
import numpy as np
from sklearn.neighbors import NearestNeighbors


def puntuacio_knn(X, k=5):
    """Distancia mitjana als k veins mes propers (sense comptar-se a si mateix).

    Utilitza un arbre KD: cost mitja ~O(n log n) en lloc d'O(n^2).
    """
    if not 1 <= k < len(X):
        raise ValueError("k ha d'estar entre 1 i n-1")
    nn = NearestNeighbors(n_neighbors=k + 1).fit(X)
    dist, _ = nn.kneighbors(X)
    return dist[:, 1:].mean(axis=1)      # la columna 0 es el propi punt


def top_anomalies(puntuacions, m=5):
    """Retorna els indexos de les m puntuacions mes altes (de major a menor)."""
    return np.argsort(puntuacions)[::-1][:m]
\end{lstlisting}

\noindent\emph{Nota:} si hi ha punts duplicats exactes, el primer veí retornat
per \texttt{kneighbors} pot no ser el propi punt; l'efecte és negligible en
aquest conjunt de dades, però es comprova a les proves de validació.

\subsection{Eficiència i ús de biblioteques}
La versió de partida reimplementa en Python pur operacions que ofereixen
NumPy i \textit{scikit-learn}. Per a un catàleg de $n=\,$\pend{valor} esdeveniments, el
càlcul $O(n^2)$ implica unes $n^2\approx\,$\pend{valor} distàncies calculades amb
bucles de Python. La taula~\ref{tab:rendiment} compara temps mesurats a la
mateixa màquina amb \texttt{time.perf\_counter} (mitjana de 5 execucions).

\begin{table}[htbp]
\centering
\small
\begin{tabular}{@{}lrrr@{}}
\toprule
\textbf{Etapa} & \textbf{Python pur (s)} & \textbf{NumPy/sklearn (s)} & \textbf{Acceleració} \\
\midrule
\textit{K-means} ($K=3$) & \pend{t1} & \pend{t2} & \pend{t1/t2}$\times$ \\
Puntuació $s_5$ & \pend{t3} & \pend{t4} & \pend{t3/t4}$\times$ \\
\bottomrule
\end{tabular}
\caption{Temps d'execució de la versió de partida i de la versió optimitzada.}
\label{tab:rendiment}
\end{table}

\subsection{Proves automàtiques}\label{sec:proves}
Es verifiquen casos normals, límit i erronis amb \texttt{pytest}
(codi~\ref{lst:tests}). El resultat de l'execució és
\pend{n} proves superades de \pend{n} (\texttt{pytest -q}).

\begin{lstlisting}[caption={Proves automàtiques (\texttt{tests/test\_sismes.py}).},label={lst:tests}]
import numpy as np
import pytest
from sismes.dades import carregar_csv
from sismes.preprocessament import estandarditzar, cartesianes_esfera
from sismes.clustering import kmeans
from sismes.anomalies import puntuacio_knn, top_anomalies


# ---- Casos normals ---------------------------------------------------------
def test_estandarditzar_mitjana_zero_desviacio_un():
    X = np.array([[1., 10.], [2., 20.], [3., 30.]])
    Z, _, _ = estandarditzar(X)
    assert np.allclose(Z.mean(axis=0), 0.0)
    assert np.allclose(Z.std(axis=0), 1.0)


def test_kmeans_separa_dos_blocs():
    rng = np.random.default_rng(0)
    X = np.vstack([rng.normal(0, .1, (50, 2)), rng.normal(10, .1, (50, 2))])
    etiq, _, _ = kmeans(X, k=2, llavor=1)
    assert len(set(etiq[:50])) == 1 and len(set(etiq[50:])) == 1
    assert etiq[0] != etiq[99]


def test_knn_detecta_outlier_sintetic():
    rng = np.random.default_rng(0)
    X = np.vstack([rng.normal(0, 1, (200, 2)), [[50., 50.]]])
    assert top_anomalies(puntuacio_knn(X, k=5), m=1)[0] == 200


def test_esfera_longitud_periodica():
    p1 = cartesianes_esfera(np.array([0.]), np.array([179.]))
    p2 = cartesianes_esfera(np.array([0.]), np.array([-179.]))
    # propers a l'esfera tot i diferir 358 graus en longitud
    assert np.linalg.norm(p1 - p2) < 300      # km (aprox. 222 km)


# ---- Casos limit -----------------------------------------------------------
def test_estandarditzar_columna_constant():
    X = np.array([[5., 1.], [5., 2.], [5., 3.]])
    Z, _, _ = estandarditzar(X)
    assert np.all(np.isfinite(Z))


def test_kmeans_k_igual_a_un():
    X = np.random.default_rng(0).normal(size=(20, 3))
    etiq, centres, _ = kmeans(X, k=1)
    assert np.all(etiq == 0)
    assert np.allclose(centres[0], X.mean(axis=0))


# ---- Casos erronis -----------------------------------------------------------
@pytest.mark.parametrize("k", [0, 6])
def test_kmeans_k_invalid(k):
    with pytest.raises(ValueError):
        kmeans(np.zeros((5, 2)), k=k)


def test_knn_k_invalid():
    with pytest.raises(ValueError):
        puntuacio_knn(np.zeros((5, 2)), k=5)


def test_carregar_csv_ignora_files_incompletes(tmp_path):
    f = tmp_path / "t.csv"
    f.write_text("mag,depth,latitude,longitude,place\n5.1,10,1,2,A\n,,,,B\n")
    X, ids = carregar_csv(str(f))
    assert X.shape == (1, 4) and ids == ["A"]


def test_carregar_csv_inexistent():
    with pytest.raises(FileNotFoundError):
        carregar_csv("no_existeix.csv")


def test_carregar_csv_sense_files_valides(tmp_path):
    f = tmp_path / "buit.csv"
    f.write_text("mag,depth,latitude,longitude,place\n")
    with pytest.raises(ValueError):
        carregar_csv(str(f))
\end{lstlisting}

\subsection{Reproducibilitat i coherència amb l'informe}\label{sec:reproducibilitat}
Per garantir que els resultats d'aquest informe es reprodueixen exactament:
\begin{itemize}[nosep]
  \item \textbf{\textit{README}:} descriu l'objectiu, l'estructura, com crear
        un entorn virtual, instal·lar dependències, executar
        \texttt{python main.py} i \texttt{pytest}, i com descarregar un nou
        catàleg.
  \item \textbf{\texttt{requirements.txt} amb versions fixades:}
        \texttt{numpy==\pend{versió}}, \texttt{scikit-learn==\pend{versió}},
        \texttt{matplotlib==\pend{versió}}, \texttt{pytest==\pend{versió}}
        (obtingudes amb \texttt{pip freeze}); es documenta la versió de
        Python \pend{versió}.
  \item \textbf{Dades d'exemple:} el fitxer \texttt{earthquakes.csv} utilitzat
        a l'informe s'inclou al repositori. Aquest és el que s'ha d'emprar per
        obtenir exactament les xifres de les taules; un catàleg més recent
        donarà valors diferents.
  \item \textbf{Llavor fixa:} \texttt{llavor=42} en totes les etapes
        aleatòries.
  \item \textbf{Instruccions d'execució:}
\end{itemize}
\begin{lstlisting}[language={},numbers=none]
python -m venv .venv && source .venv/bin/activate
pip install -r requirements.txt
python main.py          # reprodueix les taules i figures de l'informe
pytest -q               # executa les proves automatiques
\end{lstlisting}

% ==============================================================================
\section{Resultats i anàlisi}\label{sec:resultats}
% ==============================================================================
\noindent\emph{Nota per a l'autor/a: les xifres marcades en vermell s'han
d'omplir amb la sortida real de \texttt{main.py}; no s'han inventat valors
perquè l'informe ha de coincidir exactament amb l'execució.}

\subsection{Descripció del conjunt de dades}
\begin{table}[htbp]
\centering
\small
\begin{tabular}{@{}lr@{}}
\toprule
\textbf{Concepte} & \textbf{Valor} \\
\midrule
Finestra temporal & \pend{inici} -- \pend{fi} \\
Registres llegits & \pend{n$_0$} \\
Registres descartats (invàlids) & \pend{n$_0-n$} \\
Esdeveniments analitzats ($n$) & \pend{n} \\
Magnitud: mínim / mediana / màxim & \pend{.} / \pend{.} / \pend{.} \\
Profunditat (km): mínim / mediana / màxim & \pend{.} / \pend{.} / \pend{.} \\
\bottomrule
\end{tabular}
\caption{Resum del conjunt de dades.}\label{tab:dades}
\end{table}

\figuraopt{fig_distribucions.pdf}{Distribució de magnituds i profunditats i
histograma de $\log_{10}N(\ge M)$ per comprovar la llei de
Gutenberg--Richter i la completesa del catàleg.}{fig:distribucions}

La figura~\ref{fig:distribucions} permet comprovar si el catàleg segueix la
llei de Gutenberg--Richter per sobre d'una magnitud de completesa
$M_c=\,$\pend{valor}. Això es fa abans del \textit{clustering} perquè la
barreja d'esdeveniments petits (detectats només en zones ben instrumentades)
i grans pot determinar l'estructura de l'agrupació.

\subsection{Selecció del nombre de grups}
\figuraopt{fig_seleccio_k.pdf}{Inèrcia (colze), silueta mitjana i índex de
Davies--Bouldin per a $K=2,\dots,10$ (representacions A i B).}{fig:seleccio_k}

\begin{table}[htbp]
\centering
\small
\begin{tabular}{@{}rrrrrr@{}}
\toprule
$K$ & Inèrcia (A) & Silueta (A) & DB (A) & Silueta (B) & DB (B) \\
\midrule
2 & \pend{.} & \pend{.} & \pend{.} & \pend{.} & \pend{.} \\
3 & \pend{.} & \pend{.} & \pend{.} & \pend{.} & \pend{.} \\
4 & \pend{.} & \pend{.} & \pend{.} & \pend{.} & \pend{.} \\
\multicolumn{6}{c}{\dots}\\
\bottomrule
\end{tabular}
\caption{Mètriques internes per a diferents valors de $K$.}\label{tab:selecciok}
\end{table}

Segons la figura~\ref{fig:seleccio_k} i la taula~\ref{tab:selecciok}, el valor
de $K$ amb millor silueta és $K^\ast=\,$\pend{valor}, mentre que $K=3$
(el valor del codi de partida) obté una silueta de \pend{.}. \pend{Comentar si
se supera el llindar $\bar s\ge0{,}25$ de l'objectiu O2 i, si no, dir-ho
explícitament.}

\subsection{Resultats del \textit{clustering} ($K=3$)}
\begin{table}[htbp]
\centering
\small
\begin{tabular}{@{}lrrrr>{\raggedright\arraybackslash}p{4.2cm}@{}}
\toprule
\textbf{Grup} & \textbf{$n_k$} & \textbf{Magnitud mitjana} & \textbf{Profunditat mitjana (km)} & \textbf{\% del total} & \textbf{Interpretació} \\
\midrule
1 & \pend{.} & \pend{.} & \pend{.} & \pend{.} & \pend{p.~ex. sismicitat superficial} \\
2 & \pend{.} & \pend{.} & \pend{.} & \pend{.} & \pend{.} \\
3 & \pend{.} & \pend{.} & \pend{.} & \pend{.} & \pend{.} \\
\bottomrule
\end{tabular}
\caption{Característiques dels grups trobats amb $K=3$ (valors en unitats
originals).}\label{tab:clusters}
\end{table}

\figuraopt{fig_clusters.pdf}{Magnitud enfront de profunditat (eix invertit),
coloreu per grup. Els quadrats vermells marquen les 5 anomalies principals.}{fig:clusters}

\figuraopt{fig_mapa.pdf}{Mapa dels epicentres coloreats segons el grup
\textit{K-means}.}{fig:mapa}

\paragraph{Interpretació crítica.} La figura~\ref{fig:clusters} només mostra
dues de les quatre dimensions utilitzades. Cal advertir que, en incloure
latitud i longitud en l'espai de característiques (variant A), els grups no
corresponen exclusivament a règims de profunditat: una part de la separació és
\emph{geogràfica}. Per això el nom \emph{superficial/intermedi/profund} només
és apropiat si la taula de contingència
(taula~\ref{tab:contingencia}) ho confirma. \pend{Descriure què s'observa
realment: quins grups domina cada regió, si apareix un grup dominat per una
seqüència de rèpliques, etc.}

\subsection{Validació del \textit{clustering}}\label{sec:validacio}
\begin{table}[htbp]
\centering
\small
\begin{tabular}{@{}lrrr@{}}
\toprule
 & Superficial ($<70$ km) & Intermedi (70--300 km) & Profund ($>300$ km) \\
\midrule
Grup 1 & \pend{.} & \pend{.} & \pend{.} \\
Grup 2 & \pend{.} & \pend{.} & \pend{.} \\
Grup 3 & \pend{.} & \pend{.} & \pend{.} \\
\bottomrule
\end{tabular}
\caption{Taula de contingència entre els grups i la classificació de
referència per profunditat.}\label{tab:contingencia}
\end{table}

\begin{table}[htbp]
\centering
\small
\begin{tabular}{@{}lrr@{}}
\toprule
\textbf{Mètrica de validació} & \textbf{Variant A} & \textbf{Variant B} \\
\midrule
Silueta mitjana ($K=3$) & \pend{.} & \pend{.} \\
Davies--Bouldin ($K=3$) & \pend{.} & \pend{.} \\
ARI mitjà entre 30 llavors & \pend{.} & \pend{.} \\
ARI amb la classificació per profunditat & \pend{.} & \pend{.} \\
NMI amb la classificació per profunditat & \pend{.} & \pend{.} \\
\bottomrule
\end{tabular}
\caption{Validació interna, estabilitat i acord amb la referència.}
\label{tab:validacio}
\end{table}

La taula~\ref{tab:validacio} resumeix el compliment dels objectius O2 i O3.
\pend{Discutir: si l'estabilitat és $\ge0{,}90$ i si l'ARI amb la referència és
baix, explicar-ho (p.~ex. perquè la posició geogràfica domina la
distància).} També cal comentar la diferència entre les variants A i B.

\subsection{Resultats de la detecció d'anomalies}
\begin{table}[htbp]
\centering
\small
\begin{tabular}{@{}rp{4.8cm}rrrrr@{}}
\toprule
\textbf{\#} & \textbf{Esdeveniment} & $M$ & $h$ (km) & Lat. & Lon. & $s_5$ \\
\midrule
1 & \pend{lloc} & \pend{.} & \pend{.} & \pend{.} & \pend{.} & \pend{.} \\
2 & \pend{lloc} & \pend{.} & \pend{.} & \pend{.} & \pend{.} & \pend{.} \\
3 & \pend{lloc} & \pend{.} & \pend{.} & \pend{.} & \pend{.} & \pend{.} \\
4 & \pend{lloc} & \pend{.} & \pend{.} & \pend{.} & \pend{.} & \pend{.} \\
5 & \pend{lloc} & \pend{.} & \pend{.} & \pend{.} & \pend{.} & \pend{.} \\
\bottomrule
\end{tabular}
\caption{Els cinc esdeveniments amb la puntuació d'anomalia $s_5$ més alta.}
\label{tab:anomalies}
\end{table}

\figuraopt{fig_anomalies_score.pdf}{Distribució de la puntuació $s_5$ (escala
logarítmica) amb el llindar dels cinc esdeveniments principals.}{fig:scores}

\paragraph{Anàlisi.} Per a cada esdeveniment de la taula~\ref{tab:anomalies}
s'indica quina variable en justifica la puntuació (valors $z$ més alts en
magnitud, profunditat o posició). \pend{Interpretar: p.~ex. si els anòmals són
sobretot magnituds molt altes, sismes molt profunds o epicentres aïllats; i
comparar amb els esdeveniments significatius publicats per l'USGS.}

\subsection{Validació de la detecció d'anomalies}
\begin{table}[htbp]
\centering
\small
\begin{tabular}{@{}lrr@{}}
\toprule
\textbf{Prova} & \textbf{$s_5$ (kNN)} & \textbf{Observació} \\
\midrule
Recuperació d'anomalies sintètiques (top-$m_s$) & \pend{.\%} &
  Objectiu O4: $\ge 90\,\%$ \\
Jaccard top-5 amb \textit{LOF} & \pend{.} & \\
Jaccard top-5 amb \textit{Isolation Forest} & \pend{.} & \\
Variació del top-5 per a $k\in\{3,5,10,20\}$ & \pend{.} & Sensibilitat a $k$ \\
\bottomrule
\end{tabular}
\caption{Validació de la puntuació d'anomalia.}\label{tab:validacio_anom}
\end{table}

\subsection{Discussió i limitacions}
\begin{itemize}[nosep]
  \item \textbf{Anòmal estadístic $\neq$ anòmal físic.} Un terratrèmol gran
        és anòmal en l'espai de característiques perquè és rar (llei de
        Gutenberg--Richter), no perquè sigui anormal des del punt de vista
        geofísic. Un mètode no supervisat no pot fer aquesta distinció sense
        coneixement addicional.
  \item \textbf{Hipòtesis de \textit{K-means}.} Els règims tectònics (fosses,
        dorsals) són estructures allargades i \textit{K-means} tendeix a
        tallar-les en grups compactes.
  \item \textbf{Geometria de les coordenades.} La variant A ignora la
        periodicitat de la longitud i la curvatura de la Terra; la variant B
        ho corregeix parcialment.
  \item \textbf{Completesa del catàleg i seqüències de rèpliques.} Una
        seqüència intensa en el mes d'estudi pot generar un grup només per
        proximitat espacial i temporal, sense significat tectònic propi.
  \item \textbf{Finestra curta.} Un sol mes no permet inferir taxes de
        sismicitat estables; els resultats són descriptius, no predictius.
  \item \textbf{Ponderació implícita.} Després de \textit{z-score}, totes les
        variables pesen igual; això és una decisió, no una veritat física.
\end{itemize}

% ==============================================================================
\section{Conclusions i línies de millora}\label{sec:conclusions}
% ==============================================================================
\paragraph{Resposta als objectius.}
\begin{itemize}[nosep]
  \item \textbf{O1:} s'ha construït un procés de càrrega i neteja que
        documenta els registres descartats (\pend{n$_0-n$} de \pend{n$_0$}).
  \item \textbf{O2:} el valor de $K$ escollit és \pend{K}, amb silueta mitjana
        \pend{.} i ARI d'estabilitat \pend{.}; l'objectiu \pend{s'ha assolit /
        no s'ha assolit}.
  \item \textbf{O3:} l'acord amb la classificació per profunditat és
        $\mathrm{ARI}=\,$\pend{.} i $\mathrm{NMI}=\,$\pend{.}; \pend{interpretar}.
  \item \textbf{O4:} la puntuació $s_5$ recupera el \pend{.\%} de les anomalies
        sintètiques i els cinc esdeveniments principals són
        \pend{resum}.
  \item \textbf{O5:} totes les proves automàtiques passen i la versió
        optimitzada és \pend{.}$\times$ més ràpida que l'original.
\end{itemize}

\paragraph{Conclusió general.} \textit{K-means} i la distància als veïns són
eines útils com a primera exploració d'un catàleg sísmic: ofereixen un resum
interpretable i un ordenament d'esdeveniments per excepcionalitat. Tanmateix,
la validació mostra que els resultats han de llegir-se com a descripció
estadística i no com a classificació tectònica, perquè (i) les variables
geogràfiques condicionen l'agrupació, (ii) les anomalies detectades reflecteixen
principalment raresa estadística i (iii) el període d'un mes pot estar dominat
per seqüències de rèpliques.

\paragraph{Línies de millora concretes.}
\begin{enumerate}[nosep]
  \item Agrupar només amb magnitud i profunditat (i fer l'anàlisi geogràfica
        a part) o fer servir la distància de cercle màxim amb \textit{DBSCAN}
        o HDBSCAN.
  \item Incorporar el temps d'origen i aplicar agrupació espai-temps per
        separar rèpliques de la sismicitat de fons \cite{zaliapin2013}.
  \item Utilitzar \textit{k-means++}, múltiples reinicis i mescles gaussianes
        per grups el·líptics.
  \item Ampliar la finestra temporal i filtrar per magnitud de completesa
        $M_c$.
  \item Comparar sistemàticament amb \textit{LOF} i \textit{Isolation Forest}
        i, quan existeixin etiquetes (esdeveniments significatius), passar a
        una avaluació amb precisió i \textit{recall}.
\end{enumerate}

% ==============================================================================
\begin{thebibliography}{99}\addcontentsline{toc}{section}{Referències}

\bibitem{arthur2007} Arthur, D.; Vassilvitskii, S. (2007).
\textit{k-means++: The advantages of careful seeding}.
Proceedings of the 18th Annual ACM-SIAM Symposium on Discrete Algorithms
(SODA), pp.~1027--1035.

\bibitem{benioff1949} Benioff, H. (1949). Seismic evidence for the fault origin
of oceanic deeps. \textit{Bulletin of the Geological Society of America},
60(12), 1837--1856.

\bibitem{bird2003} Bird, P. (2003). An updated digital model of plate
boundaries. \textit{Geochemistry, Geophysics, Geosystems}, 4(3), 1027.

\bibitem{breunig2000} Breunig, M. M.; Kriegel, H.-P.; Ng, R. T.; Sander, J.
(2000). LOF: Identifying density-based local outliers.
\textit{Proceedings of the ACM SIGMOD International Conference on Management of
Data}, pp.~93--104.

\bibitem{chandola2009} Chandola, V.; Banerjee, A.; Kumar, V. (2009). Anomaly
detection: A survey. \textit{ACM Computing Surveys}, 41(3), 15:1--15:58.

\bibitem{davies1979} Davies, D. L.; Bouldin, D. W. (1979). A cluster
separation measure. \textit{IEEE Transactions on Pattern Analysis and Machine
Intelligence}, 1(2), 224--227.

\bibitem{ester1996} Ester, M.; Kriegel, H.-P.; Sander, J.; Xu, X. (1996). A
density-based algorithm for discovering clusters in large spatial databases
with noise. \textit{Proceedings of the 2nd International Conference on
Knowledge Discovery and Data Mining (KDD)}, pp.~226--231.

\bibitem{gutenberg1944} Gutenberg, B.; Richter, C. F. (1944). Frequency of
earthquakes in California. \textit{Bulletin of the Seismological Society of
America}, 34(4), 185--188.

\bibitem{hanks1979} Hanks, T. C.; Kanamori, H. (1979). A moment magnitude
scale. \textit{Journal of Geophysical Research}, 84(B5), 2348--2350.

\bibitem{hubert1985} Hubert, L.; Arabie, P. (1985). Comparing partitions.
\textit{Journal of Classification}, 2(1), 193--218.

\bibitem{kaufman1990} Kaufman, L.; Rousseeuw, P. J. (1990).
\textit{Finding Groups in Data: An Introduction to Cluster Analysis}.
Wiley.

\bibitem{liu2008} Liu, F. T.; Ting, K. M.; Zhou, Z.-H. (2008). Isolation
forest. \textit{Proceedings of the 8th IEEE International Conference on Data
Mining (ICDM)}, pp.~413--422.

\bibitem{lloyd1982} Lloyd, S. P. (1982). Least squares quantization in PCM.
\textit{IEEE Transactions on Information Theory}, 28(2), 129--137.

\bibitem{macqueen1967} MacQueen, J. (1967). Some methods for classification
and analysis of multivariate observations. \textit{Proceedings of the 5th
Berkeley Symposium on Mathematical Statistics and Probability}, vol.~1,
pp.~281--297.

\bibitem{pedregosa2011} Pedregosa, F. et al. (2011). Scikit-learn: Machine
learning in Python. \textit{Journal of Machine Learning Research}, 12,
2825--2830.

\bibitem{ramaswamy2000} Ramaswamy, S.; Rastogi, R.; Shim, K. (2000). Efficient
algorithms for mining outliers from large data sets. \textit{Proceedings of the
ACM SIGMOD International Conference on Management of Data}, pp.~427--438.

\bibitem{rousseeuw1987} Rousseeuw, P. J. (1987). Silhouettes: A graphical aid
to the interpretation and validation of cluster analysis. \textit{Journal of
Computational and Applied Mathematics}, 20, 53--65.

\bibitem{sinnott1984} Sinnott, R. W. (1984). Virtues of the haversine.
\textit{Sky and Telescope}, 68(2), 159.

\bibitem{usgs} U.S. Geological Survey. \textit{Earthquake Catalog (ComCat)} i
canals d'esdeveniments. \url{https://earthquake.usgs.gov/earthquakes/search/}.
Consultat el \pend{data de consulta}.

\bibitem{zaliapin2013} Zaliapin, I.; Ben-Zion, Y. (2013). Earthquake clusters in
southern California I: Identification and stability. \textit{Journal of
Geophysical Research: Solid Earth}, 118(6), 2847--2864.

\end{thebibliography}

\end{document}
